\section{Project Strategy \& Scope}

\subsection{Mission Statement}
The primary objective of \textbf{Project Pienza} is to deconstruct, model, and optimize the decision-making process of a rational agent within the ride-hailing ecosystem. The environment is formally defined as a \textit{sequential, non-cooperative game} governed by intrinsic information asymmetry.

Unlike traditional "top-down" studies relying on aggregated data, Pienza employs a \textbf{"Bottom-Up"} methodology grounded in the operational reality of Mexico City. Through the longitudinal quantification of a single Expert Agent ($N=1$), this research models the micro-economic physics of the Gig Economy.

\subsection{Scope \& Constraints}
The project explicitly rejects the naive premise of reverse-engineering the proprietary pricing algorithms of ride-hailing platforms. Modeling a dynamic, global algorithm using a local boutique dataset is stochastically infeasible.

Instead, Pienza reorients the analytical lens—shifting away from the fare prediction models common in public repositories—to focus on the sole variable under absolute control: \textbf{The Agent's Decision}. By pivoting to Behavioral Economics and Game Theory, the project transforms an intractable prediction problem into a solvable policy optimization endeavor.


\subsection{Core Framework}
The project is structured across three interdependent modules covering the data science lifecycle:

\begin{enumerate} 
    \item \textbf{Data Ingestion \& Infrastructure:} 
    A custom ingestion pipeline was built to transform unstructured operational data into a structured relational schema. This involved the deployment of OCR-based capture and geospatial triangulation logic to digitize the agent's real-time decision environment.
    
    \item \textbf{Decision Stream Analysis (The Negative Class):} 
    Traditional market analyses focus exclusively on completed transactions, resulting in \textit{Survivorship Bias}. Project Pienza captures the \textbf{Total Offer Stream}, recording thousands of rejected offers. This granular visibility into the decision boundary allows for the quantification of:
    \begin{itemize}
        \item \textbf{Excess Return Coefficient:} Isolation of marginal gains generated by the agent's strategy relative to the market baseline.
        \item \textbf{Optimal Stopping:} Modeling the threshold where the cost of search (deadhead time) exceeds the probable value of future offers.
        \item \textbf{Risk Transfer:} Measuring the variance between algorithmic duration forecasts and operational reality.
    \end{itemize}
    
    \item \textbf{Policy Cloning \& Simulation:}
    The final analytical objective is the transition from descriptive analysis to predictive modeling. By training \textit{Imitation Learning} models using XGBoost, Pienza generates a digital clone of the agent's heuristics. This model enables high-fidelity market simulations and counterfactual testing using synthetic data generated via GAN architectures.
\end{enumerate}


\subsection{Initial Hypothesis: The Payout Spread}
Research began with a longitudinal field observation regarding fare elasticity. The hypothesis stated that the system applied an implicit efficiency penalty, observed as a discrepancy between the promised fare and realized earnings. A preliminary pilot study ($N \approx 150$, July-August cohort) revealed a \textbf{Payout Rate (Spread)} between $75\%$ and $85\%$ of the Base Fare. The distribution was relatively uniform and followed \textit{session-based moving averages}, suggesting intra-block temporal stability driven by market conditions.

\subsection{Model Performance and Data Scarcity}
An initial model correlated this Spread with the \textbf{Time Delta} (the discrepancy between estimated and realized trip time). The goal was to validate whether time savings generated a fare penalty.

\textbf{Key Finding:} A benchmark comparison showed that Simple Linear Regression outperformed base Machine Learning models. No significant non-linear relationship was identified; the correlation remained weak. While basic statistics resolved the regression problem for the pilot sample, scaling this to a robust ML model would require thousands of \textit{completed} trips, representing an infeasible operational time cost for marginal analytical return.

\subsection{Strategic Pivot to Classification}
Based on logistical constraints and the low predictive value of the initial regression approach, the research objective was redefined:

\begin{itemize}
    \item \textbf{From:} Price variance prediction (Regression), which required scarce completed trip data.
    \item \textbf{To:} Behavior modeling (Classification), which utilized the total offer stream.
\end{itemize}

This shift resolved the data scarcity issue by enabling the inclusion of the \textbf{Negative Class} (rejected offers). By modeling decision policy rather than price, the project could apply high-dimensional Machine Learning (XGBoost) to capture the agent's non-linear decision boundary.

\subsection{Project Roadmap}
The project followed a non-linear development lifecycle where feature engineering evolved organically in response to the analytical findings of each stage:

\begin{description}[style=nextline, leftmargin=0.5cm, font=\bfseries]
    \item[\textcolor{red}{Phase 1: Acquisition \& Ground Truth (Aug 22 -- Oct 1, 2025)}]
    \textcolor{red}{\textbf{[PENDING UPDATE]} High-fidelity manual capture of operational data via screenshots and the \textit{GTS Webapp}. Establishment of the target variable and the \textbf{Bronze Palette} (contextual ride attributes integrated directly into the core dataset).}

    \item[\textcolor{red}{Phase 2: Data Engineering \& Architecture (Oct 2 -- Nov 20, 2025)}]
    \textcolor{red}{\textbf{[PENDING UPDATE]} Transition to a relational database system (\texttt{pienza.db}) and Star Schema design. Implementation of automated OCR pipelines and normalization protocols. Creation of the \textbf{Silver Palette} (\texttt{engineered\_features} table).}

    \item[\textcolor{red}{Phase 3: Exploratory Analyses and Causal Inference (Nov 21 -- Dec 3, 2025)}]
    \textcolor{red}{\textbf{[PENDING UPDATE]} Forensic audit of market friction and validation of the ``Risk Transfer Thesis.'' Execution of dimensionality reduction (Lasso, PCA) to identify primary predictive signals.}

    \item[\textcolor{red}{Phase 4: Unsupervised Learning \& Geo-Remediation (Dec 4, 2025 -- Jan 2, 2026)}]
    \textcolor{red}{\textbf{[PENDING UPDATE]} Application of K-Means and HDBSCAN for zone discovery, combined with surgical coordinate cleaning using hand-crafted heuristic polygons. Generation of the \textbf{Gold Palette} (geo-semantic attributes).}

    \item[\textcolor{red}{Phase 5: Supervised Learning (Jan 3 -- Jan 12, 2026)}]
    \textcolor{red}{\textbf{[PENDING UPDATE]} Model evaluation tournament (Naive Bayes vs. LogReg vs. XGBoost). Implementation of the \textit{Cognitive Cascade} hierarchical architecture for the champion model.}

    \item[\textcolor{red}{Phase 6: Generative Simulation (Jan 13, 2026 -- Ongoing)}]
    \textcolor{red}{\textbf{[PENDING UPDATE]} Infrastructure migration to Google BigQuery. Deployment of GANs for synthetic manifold generation and policy stress-testing under simulated market conditions.}
\end{description}


\subsection{A Note on System Architecture \& Evolution}
The rigorous architecture of the data ecosystem was established as a foundational objective, carrying equal weight to the subsequent analytical and predictive modeling phases. The infrastructure was not static; it evolved through four distinct operational states to match the complexity of the inquiry:

\begin{itemize}
    \item \textbf{Phase 1 (Acquisition Staging):} Data ingestion from Engines 1 \& 2 was routed into \textit{Google Sheets}, supported by ad-hoc Python and Google Apps Scripts for initial validation.
    \item \textbf{Phase 2 (The Engineering Core):} The transition from flat spreadsheets to a relational architecture. Execution of the idempotent ETL pipeline to forge the definitive \textbf{SQLite} asset (\texttt{pienza.db}).
    \item \textbf{Phases 3--5 (The Analytical Loop):} A file-based cloud architecture where Python Notebooks (EDA, Clustering, Supervised Learning) interacted with \texttt{pienza.db} as the immutable Single Source of Truth.
    \item \textbf{Phase 6 (Cloud Scaling):} Migration to \textbf{Google BigQuery} to support Generative AI workflows.
    \begin{itemize}
        \item \textbf{\texttt{pienza\_mini}:} A high-fidelity cloud replica of the original ground-truth database, preserving the identical relational logic.
        \item \textbf{\texttt{pienza\_manifold}:} A 1,000,000-row synthetic dataset generated via GANs. The Parquet artifacts are persisted on Google Drive and queried through BigQuery's external table interface to ensure zero-latency analysis.
    \end{itemize}
\end{itemize}
